\section{Regime Robustness Study}
\label{sec:robustness-plan}

The primary evaluation uses one station layout, one CBRN-like hazard placement,
one mixed-familiarity population, and a fixed evacuation horizon. That is enough
to identify the information-safety tradeoff, but not enough to claim that the
observed policy ordering generalizes. The next empirical package therefore
treats external validity as the object of study.

The robustness study asks where each communication policy becomes safe,
inefficient, or harmful. It varies three hazard regimes, three population
familiarity regimes, and five representative communication policies over 20
seeds, yielding 900 runs. The policy set is intentionally smaller than the main
study: no intervention, static beacons, global broadcast, entropy-targeted
communication, and bottleneck-avoidance communication. These policies span the
main theoretical contrast between no control, local fixed control, global
high-reach control, uncertainty-driven adaptive control, and capacity-aware
adaptive control.

\begin{table*}[!t]
\centering
\caption{Planned regime robustness matrix. The full factorial design contains
$3 \times 3 \times 5 \times 20 = 900$ seeded runs.}
\label{tab:robustness-matrix}
\begin{tabular}{p{0.18\textwidth}p{0.34\textwidth}p{0.38\textwidth}}
\toprule
Axis & Levels & Purpose \\
\midrule
Hazard regime & Low, medium, high severity/spread & Tests whether information
control changes under stronger physical threat. \\
Population familiarity & Low, mixed, high exit familiarity & Tests whether
communication helps unfamiliar crowds and whether familiar crowds synchronize. \\
Communication policy & None, static, global, entropy, bottleneck & Tests the
core safety-control frontier without diluting the study with every policy. \\
\bottomrule
\end{tabular}
\end{table*}

This study is not designed to maximize a single score. Its role is to estimate
an operating envelope: when local communication remains efficient, when global
broadcast becomes wasteful or risky, when entropy targeting needs budget
control, and when bottleneck-aware messaging is most valuable. If the main
claim is correct, the robustness grid should show that the sign and magnitude
of information-safety efficiency depend on hazard pressure and prior
familiarity, while broad reach remains an unreliable proxy for safe control.

\subsection{Interpretation Protocol}
\label{subsec:robustness-interpretation}

The robustness grid will be interpreted as a claim filter rather than as a
leaderboard. A policy will be treated as \emph{globally robust} only if its
information-safety efficiency remains positive and competitive across most
hazard--familiarity cells without producing consistently high harmful
convergence. A policy will be treated as \emph{regime dependent} if it performs
well only under particular hazard severity or familiarity conditions. A policy
will be treated as \emph{unsafe or inefficient under this model} if it improves
belief state while repeatedly increasing exposure, queue pressure, or exit
concentration.

This protocol matters because the paper's main contribution is a measurement
framework, not a universal evacuation prescription. For example, static beacons
may remain the strongest conservative controller when hazards are mild and
prior familiarity is low, while bottleneck-aware messaging may become more
valuable under higher crowd synchronization or stronger hazard pressure.
Similarly, entropy targeting may be scientifically important even when it is
not operationally dominant, because it reveals when uncertainty reduction
turns into synchronized route choice. The final paper will therefore report
both policy rankings and regime-specific failure modes.

\subsection{Decision Criteria for the Final Paper}
\label{subsec:robustness-decision-criteria}

The completed robustness study will update the claim matrix in
\S\ref{sec:claims}. If static beacons or bottleneck-aware messaging retain
positive ISE across most regimes, the paper can claim that local or
capacity-aware information control is a robust baseline for this simulator
family. If global broadcast remains lower-efficiency despite broad reach, the
paper can strengthen the claim that communication reach is not an adequate
surrogate for safety. If entropy targeting performs well only in limited cells,
the paper should frame it as a diagnostic control mechanism rather than as a
deployment recommendation. If policy ordering changes sharply across the grid,
the paper should emphasize operating envelopes and avoid any single-policy
conclusion.

The study can also falsify parts of the current interpretation. If broad
broadcast becomes consistently high-efficiency under high hazard severity, the
paper must narrow the ``reach is not enough'' claim to lower-pressure or
capacity-constrained settings. If exposure-aware or entropy-targeted control
unexpectedly dominates the grid, the negative result from the primary study
should be presented as scenario-specific. These are useful outcomes: they would
make the paper stronger by showing where the information-safety framework
changes conclusions as the physical and informational regime changes.
